\documentclass[10pt,a4paper]{amsart}
\usepackage[margin=19mm]{geometry}
\usepackage{amsmath,amssymb,mathtools}
\usepackage[scr=rsfs,cal=euler]{mathalfa}
\usepackage{xfrac}
\usepackage[unicode,hidelinks,bookmarks=false]{hyperref}
\newcommand{\ie}{i.e.\ }
\newcommand{\cf}{cf.\ }
\newcommand{\resp}{respectively\ }
\newcommand{\Subsection}{\subsection}
\providecommand{\nobreakdash}{\nobreak\hbox{-}}
\setlength{\emergencystretch}{2em}
\multlinegap0pt
\usepackage{bm}
\usepackage{bbm}
\usepackage{dsfont}
\usepackage{xspace}
\usepackage{cancel}
\usepackage{mathtools}
\usepackage{enumitem}
\usepackage{amsthm}
\makeatletter
\@ifundefined{theorem}{\newtheorem{theorem}{Theorem}}{}
\@ifundefined{lemma}{\newtheorem{lemma}[theorem]{Lemma}}{}
\makeatother
\DeclareMathOperator{\grad}{grad{}}
\DeclareMathOperator{\sing}{sing \, }
\newcommand{\C}{\mathbb{C}}
\title[On Teissier's example of an equisingularity class that canno]{On Teissier's example of an equisingularity class that cannot be defined over the rationals}
\author{Adam Parusi\'nski, Laurentiu Paunescu}
\date{Source: arXiv:2501.10863v3 (CC BY 4.0)}
\begin{document}
\maketitle

\noindent\emph{Source and licence.} On Teissier's example of an equisingularity class that cannot be defined over the rationals. Version arXiv:2501.10863v3 (CC BY 4.0), distributed under the Creative Commons Attribution 4.0 International licence, as recorded in the arXiv licence metadata for this exact version and shown on the abstract page https://arxiv.org/abs/2501.10863v3. Full rights notice: licenses/arxiv-2501.10863-CC-BY-4.0.txt.\par\smallskip

\noindent\textit{Editorial scope.} Counted unit: the Lemma whose source label is \texttt{lem:nonisolates} in the article (source0.tex lines 475--485, source label \texttt{lem:nonisolates}), delivered together with the article's own complete original proof of it (source0.tex lines 489--514, one \texttt{proof} environment of 26 lines). No mathematics is written, repaired or re-derived: statement and proof are the article's own text line for line. The proof is detached from the statement in the source: the article states the result at lines 475--485 and prints its proof at lines 489--514, and the proof environment's own header names the statement, so the association is the article's own. Required context retained uncounted: none. Every definition, display and label of those lines is retained unchanged. Omitted from this package and not redistributed: 1--474, 486--488, 515--955. Nothing inside a retained excerpt is omitted or altered: every retained body is delivered exactly as the article's lines are written, so no omission receipt of any kind accompanies this package. Citations: the retained excerpts carry no citation command at all, so no bibliography entry of the article is transcribed and no reference is redistributed. Reference adaptation: none was needed and none was made; every reference inside the delivered unit resolves inside it, so no unresolved reference or citation is produced. Printed slips kept literal: no printed word, symbol, hypothesis or number of the counted statement or of its proof was altered. Layout and class: the article's own class is \texttt{amsart} with packages including amssymb, hyperref, enumitem, graphicx, tikz, tikz-cd, multicol, fontenc; the wrapper is the lane's pinned \texttt{amsart} editorial wrapper at 10pt on A4, which restates the theorem-like environments the retained text needs and adds the article's own macro definitions that the retained text needs (\texttt{\textbackslash C}, \texttt{\textbackslash grad}, \texttt{\textbackslash sing}), with extra wrapper packages (bm, bbm, dsfont, xspace, cancel, mathtools, enumitem). The delivered numbering is the unit's own. Assets: no retained span contains a figure environment, a graphics-inclusion command, a PDF-page inclusion, a table or a TikZ picture; no image file is copied into this package, the package declares no asset dependency and no image file is redistributed with it. Licence: version arXiv:2501.10863v3 of this article is distributed under the Creative Commons Attribution 4.0 International licence, as recorded in the arXiv licence metadata for this exact version and shown on the abstract page https://arxiv.org/abs/2501.10863v3. A full rights notice is delivered at licenses/arxiv-2501.10863-CC-BY-4.0.txt. Characters: every retained excerpt is pure ASCII, so the delivered engine, XeLaTeX with its default Latin Modern text font, reports no missing character.
 \begin{lemma}\label{lem:nonisolates}
Let $S\subset (\C^n,0)$ be a hypersurface, $S=f^{-1} (0)$ with $f$ reduced, and let $\ell$ be a line through the origin in $\C^n$ that is not tangent to $S_{\sing}$.  Then, for a generic linear form $\varphi: \C^ n\to \C$ and for $N$ sufficiently large,   
$g(x):= f(x) + \varphi^N(x)$ satisfies 
\begin{enumerate}[label=\roman*)] 
\item  
$(g^{-1} (0))_{\sing}= S_{\sing} \cap \varphi^{-1} (0)$; 
\item 
for a sufficiently small conical open neighbourhood $U\subset \C^n$ of $\ell^*:= \ell \setminus \{0\}$,  
the limits as $x\to 0$, $x\in U$,  of hyperplanes tangent at $x$ to the levels of $f$ and those to the levels of $g$  coincide.  
\end{enumerate}
 \end{lemma}

  \begin{proof}[Proof of Lemma \ref{lem:nonisolates}] 
Let the system of coordinates $x=(x_1, \ldots , x_n)$ be such that  
\begin{align}\label{eqn:polar}
\left \{ \frac{\partial f}{\partial x_1}= \frac{\partial f}{\partial x_2}
= \cdots =\frac{\partial f}{\partial x_{n-1} }=0 \right \} = \Gamma \cup 
S_{\sing}, %\cap S_{sing},
\end{align} 
with $\Gamma$ being of dimension $1$ (or empty).  Then $\Gamma$, a generic relative polar curve,  depends only on the projection $(x_1, \ldots , x_{n-1}) : \C^n \to \C^{n-1}$.  Therefore, after changing $x_n$ if necessary, we may assume moreover that $\{x_n=0\} \cap (\Gamma \cup \ell) = \{0\}$ and   
denote $H:=\{x_n=0\}$.  
Let $U$ be a conical neighbourhood of $\ell^*$ such that  $\overline U \cap (H\cup C_{S_{\sing},0}) = \{0\}$, $0\notin U$.  Then, by the {\L}ojasiewicz inequality, there exist a positive integer $N$ and a constant $c>0$ such that 
\begin{align}\label{eqn:loj}
\|\grad f (x)\| \ge c |x_n|^{N-2} \quad \text { for } x\in U.
\end{align}

Let $g(x):= f(x) + x_n^N$.  Then 
\begin{align}\label{eqn:gandf}
\frac{\partial g}{\partial x_i} = \frac{\partial f}{\partial x_i} \text { for } i= 1, \ldots  ,n-1 
\text { and } \frac{\partial g}{\partial x_n} = \frac{\partial f}{\partial x_n} + Nx_n^{N-1}. 
\end{align}
Therefore, because $H\cap \Gamma = \{0\}$, we have $(g^{-1} (0))_{\sing}= S_{\sing} \cap \varphi^{-1} (0)$ as needed.  By \eqref{eqn:loj} and \eqref{eqn:gandf}, $g$ on $U$ satisfies the {\L}ojasiewicz inequality similar to \eqref{eqn:loj} 
 and therefore 
$$
\left | \frac{\grad f}{\|\grad f\|} -\frac{\grad g}{\|\grad g\|} \right | \le C |x_n| .
$$
Then $\varphi (x) = x_n$ and $N$ given by \eqref{eqn:loj} guarantee  the conclusions of the lemma. 
 \end{proof}

\end{document}
